% 第一章檔案：chapter1.tex
% 這是第一章的獨立內容，可直接 \include 到主檔案

\section{定義與原則}

實證醫學（Evidence-Based Medicine, EBM）被廣泛定義為「謹慎、明確且明智地使用當前最佳證據，來做出關於個別患者照護的決策」。這一定義源自David Sackett等人於1996年的經典描述，強調實證醫學並非僅依賴傳統經驗或專家意見，而是將最佳可用外部臨床證據與個別臨床專業知識相整合。

實證醫學的核心原則包括以下三個要素：
\begin{itemize}
    \item \textbf{最佳研究證據}：來自系統性研究的臨床相關證據，例如隨機對照試驗（RCT）、系統回顧與meta-analysis。
    \item \textbf{臨床專業知識}：醫師透過臨床經驗與實踐所累積的專業判斷能力，用以有效診斷並考量患者個別情況。
    \item \textbf{患者價值與偏好}：決策必須尊重患者的權利、期望與個人情境，實現共享決策（shared decision-making）。
\end{itemize}

實證醫學的實踐通常遵循五個步驟（Sackett模型）：
\begin{enumerate}
    \item 將臨床問題轉化為可回答的結構化問題（例如使用PICO框架）。
    \item 搜尋最佳可用證據。
    \item 批判性評估證據的有效性、影響力與適用性。
    \item 將證據應用於臨床實踐，並整合臨床專業與患者價值。
    \item 評估實踐成效並持續改進。
\end{enumerate}

這些原則確保醫療決策基於科學證據，而非直覺或陳規，從而提升照護品質與患者安全。

\section{EBM的歷史發展}

實證醫學的哲學起源可追溯至19世紀中葉的巴黎醫學派，以及更早的醫學統計應用。然而，現代實證醫學的正式發展始於20世紀後期。

關鍵歷史里程碑包括：
\begin{itemize}
    \item \textbf{1972年}：英國流行病學家Archie Cochrane出版《Effectiveness and Efficiency》，批評醫療資源分配缺乏隨機對照試驗支持，奠定Cochrane Collaboration的基礎。
    \item \textbf{1980年代}：加拿大McMaster大學的David Sackett領導臨床流行病學部門，發展批判性評估方法，並出版系列文章指導醫師閱讀文獻。
    \item \textbf{1990年}：Gordon Guyatt首次使用「evidence-based medicine」一詞，描述新的醫學教學方法。
    \item \textbf{1992年}：Evidence-Based Medicine Working Group發表五步驟模型。
    \item \textbf{1996年}：Sackett等人於《BMJ》發表經典論文，明確定義EBM並澄清其非「食譜式醫學」。
    \item \textbf{1990年代後期}：成立Cochrane Collaboration與Oxford Centre for Evidence-Based Medicine，推動系統回顧與全球合作。
\end{itemize}

實證醫學從個別決策擴展至指南制定與政策，現已融入醫學教育與臨床實踐。

\section{EBM在醫療決策中的角色}

實證醫學在醫療決策中扮演核心角色，它橋接科學證據與臨床實務，確保決策基於可靠證據而非主觀判斷。這有助於：
\begin{itemize}
    \item 提升決策品質：減少無效或有害介入，促進有效治療。
    \item 資源最適化：避免不必要檢查或治療，降低醫療成本。
    \item 患者中心照護：整合患者價值，促進共享決策，提高滿意度與依從性。
    \item 終身學習：鼓勵醫師持續更新知識，應對醫學快速進展。
\end{itemize}

在臨床情境中，EBM適用於診斷、治療、預後與預防等領域。例如，面對患者問題時，醫師可使用EBM步驟快速搜尋並應用最新證據，取代過時經驗。儘管面臨時間限制或證據不足等挑戰，EBM仍是現代醫學的基石，持續推動醫療品質改進。

% 可在此新增圖表、參考文獻或習題
% \bibliographystyle{plain}
% \bibliography{references} % 若有參考文獻檔案
